\mxnewpage\section*{EV3-B -- Blanc 2}\label{EV3-B}
\addcontentsline{toc}{section}{EV3-B -- Blanc 2}
\textbf{Sujet réservé au professeur}. Programme étudié des chapitres 1 à 10.\par Durée : 2 h. Sans calculatrice. Note sur 20. 1–3 avril 2027.
\subsection*{Partie 1 -- Automatismes (6 points)}
Douze questions, \num{0,5} point chacune. Une seule réponse correcte ; aucune justification demandée. Une réponse incorrecte ou absente rapporte 0, sans retrait de points.
\begin{enumerate}
\item Calculer $4+2\times3$.\par A. $10$\quad B. $9$\quad C. $18$\quad D. $13$
\item Calculer $\frac12+\frac13$.\par A. $\frac{1}{6}$\quad B. $\frac{5}{6}$\quad C. $\frac{5}{3}$\quad D. $\frac{2}{5}$
\item Calculer $\dfrac{2^4}{2^{2}}$.\par A. $4$\quad B. $1$\quad C. $2$\quad D. $9$
\item Une valeur de 50 augmente de \pc{20}. Valeur finale ?\par A. $70$\quad B. $60$\quad C. $10$\quad D. $40$
\item Une valeur de 80 est obtenue après une baisse de \pc{20}. Valeur initiale ?\par A. $96$\quad B. $60$\quad C. $64$\quad D. $100$
\item Hausse de \pc{25} puis baisse de \pc{20} : coefficient global ?\par A. $\frac{9}{10}$\quad B. $1$\quad C. $\frac{21}{20}$\quad D. $\frac{4}{5}$
\item La solution de $2x+4=0$ vaut\par A. $2$\quad B. $8$\quad C. $-2$\quad D. $-8$
\item La pente entre $\pt{0;4}$ et $\pt{2;10}$ vaut\par A. $2$\quad B. $-3$\quad C. $3$\quad D. $6$
\item La moyenne de $4, 6, 8$ vaut\par A. $7$\quad B. $18$\quad C. $6$\quad D. $4$
\item Un secteur de $90^\circ$ représente quelle fraction du total ?\par A. $\frac{1}{2}$\quad B. $\frac{1}{3}$\quad C. $\frac{3}{4}$\quad D. $\frac{1}{4}$
\item Une probabilité vaut $\dfrac{2}{5}$. Son complément vaut\par A. $\frac{1}{5}$\quad B. $\frac{5}{2}$\quad C. $\frac{2}{5}$\quad D. $\frac{3}{5}$
\item Pour $f(x)=x^2+4$, $f(-2)$ vaut\par A. $16$\quad B. $0$\quad C. $6$\quad D. $8$
\end{enumerate}
\mxnewpage\subsection*{Partie 2 -- Exercices indépendants (14 points)}
Toute démarche doit être justifiée. Les résultats exacts sont attendus sauf indication contraire.
\mxneed{12}
\subsection*{Exercice 1 -- Quatre essais et un gain (5 points)}
Un jeu comporte quatre essais indépendants ; chaque essai réussit avec probabilité $\dfrac12$. Le gain brut $G$ vaut 8 unités monétaires si les quatre essais réussissent, et 0 sinon. Une mise pourra être fixée pour rendre le jeu équitable.
\begin{enumerate}
\item Calculer la probabilité de recevoir 8. \hfill(1 pt)
\item Calculer espérance et variance de G. \hfill(2 pts)
\item Déterminer la mise équitable ; expliquer si l’espérance est un gain possible. \hfill(2 pts)
\end{enumerate}
\mxneed{12}
\subsection*{Exercice 2 -- Une dérivée à signe constant (5 points)}
On considère la fonction $f$ définie sur $\R$ par $f(x)=(x^2+1)e^x$. On veut vérifier si une tangente horizontale correspond nécessairement à un extremum.
\begin{enumerate}
\item Calculer et factoriser $f\prime$. \hfill(1 pt)
\item Étudier la variation ; la tangente horizontale en -1 donne-t-elle un extremum ? \hfill(2 pts)
\item Écrire la tangente en 0 puis calculer $f(1)$ exactement. \hfill(2 pts)
\end{enumerate}
\mxneed{12}
\subsection*{Exercice 3 -- Projection sur une diagonale (4 points)}
Dans un repère orthonormé d’origine $O$, on considère $P\pt{2;0}$ et la droite $d$ d’équation $y=x$. On note $H$ le projeté orthogonal de $P$ sur $d$. Le vecteur unitaire demandé à la fin doit être orienté de $O$ vers $H$.
\begin{enumerate}
\item Donner un vecteur normal non nul à la droite $d$. \hfill(1 pt)
\item Déterminer les coordonnées de $H$, puis la distance $PH$. \hfill(1 pt)
\item Donner l’angle non orienté $\widehat{POH}$ en radians, puis les coordonnées du vecteur $\dfrac{\vect{OH}}{OH}$. \hfill(2 pts)
\end{enumerate}
